% === Geometry.Ch2.Prop ===

\begin{theorem}
  \label{Geometry.Ch2.Prop.P3}
  \uses{Geometry.Theory.Line, Geometry.Theory.Point, Geometry.Theory.I3}
  \lean{Geometry.Ch2.Prop.P3}
  \leanok
  p71, "For every line, there is at least one point not lying on it." 
\begin{verbatim}P3 : forall (L : Geometry.Theory.Line), Exists.{1} Geometry.Theory.Point (fun (P : Geometry.Theory.Point) => Not (Membership.mem.{0, 0} Geometry.Theory.Point Geometry.Theory.Line (Set.instMembership.{0} Geometry.Theory.Point) L P))\end{verbatim}
\end{theorem}

\begin{theorem}
  \label{Geometry.Ch2.Prop.P4}
  \uses{Geometry.Theory.Point, Geometry.Theory.Line, Geometry.Theory.Concurrent, Geometry.Ch2.Prop.P2}
  \lean{Geometry.Ch2.Prop.P4}
  \leanok
  p71. "For every point, there is at least one line not passing through it." 
\begin{verbatim}P4 : forall (P : Geometry.Theory.Point), Exists.{1} Geometry.Theory.Line (fun (L : Geometry.Theory.Line) => Not (Membership.mem.{0, 0} Geometry.Theory.Point Geometry.Theory.Line (Set.instMembership.{0} Geometry.Theory.Point) L P))\end{verbatim}
\end{theorem}

\begin{theorem}
  \label{Geometry.Ch2.Prop.P5}
  \uses{Geometry.Theory.Point, Geometry.Theory.Line, Geometry.Theory.Point.distinct_points_exist, Geometry.Theory.I1, Geometry.Ch2.Prop.P3, Geometry.Theory.Line.distinguishing_point}
  \lean{Geometry.Ch2.Prop.P5}
  \leanok
  p71. "For every point P, there are at least two distinct lines through P" 
\begin{verbatim}P5 : forall (P : Geometry.Theory.Point), Exists.{1} Geometry.Theory.Line (fun (L : Geometry.Theory.Line) => Exists.{1} Geometry.Theory.Line (fun (M : Geometry.Theory.Line) => And (Ne.{1} Geometry.Theory.Line L M) (And (Membership.mem.{0, 0} Geometry.Theory.Point Geometry.Theory.Line (Set.instMembership.{0} Geometry.Theory.Point) L P) (Membership.mem.{0, 0} Geometry.Theory.Point Geometry.Theory.Line (Set.instMembership.{0} Geometry.Theory.Point) M P))))\end{verbatim}
\end{theorem}

\begin{lemma}
  \label{test_lemma}\lean{Geometry.Theory.Betweenness.abc_imp_collinear}\leanok
  test
\end{lemma}

\begin{theorem}
  \label{Geometry.Ch2.Prop.P1}
  \leanok
  \lean{Geometry.Ch2.Prop.P1}
  \uses{Geometry.Theory.Line, Geometry.Theory.Parallel, Geometry.Theory.Point, Geometry.Theory.Intersects, Geometry.Ch2.Prop.P1.direct}
  A corrolary of the main theorem that is more useful since it uses the syntax directly. 
\begin{verbatim}P1 : forall {L : Geometry.Theory.Line} {M : Geometry.Theory.Line}, (Ne.{1} Geometry.Theory.Line L M) -> (Not (Geometry.Theory.Parallel L M)) -> (ExistsUnique.{1} Geometry.Theory.Point (fun (X : Geometry.Theory.Point) => Geometry.Theory.Intersects L M X))\end{verbatim}
\end{theorem}

\begin{theorem}
  \label{Geometry.Ch2.Prop.P2}
  \uses{Geometry.Theory.Point, Geometry.Theory.Line, Geometry.Theory.Concurrent, Geometry.Theory.I3, Geometry.Theory.I1, Geometry.Theory.Line.concurrence_of_three_lines_is_unique, Geometry.Theory.Parallel, Geometry.Theory.Line.intersecting_lines_are_not_parallel, Geometry.Ch2.Prop.P1.direct}
  \lean{Geometry.Ch2.Prop.P2}
  \leanok
  p71, "There exist three distinct lines that are not concurrent." 
\begin{verbatim}P2 : Exists.{1} Geometry.Theory.Line (fun (L : Geometry.Theory.Line) => Exists.{1} Geometry.Theory.Line (fun (M : Geometry.Theory.Line) => Exists.{1} Geometry.Theory.Line (fun (N : Geometry.Theory.Line) => And (And (Ne.{1} Geometry.Theory.Line L M) (And (Ne.{1} Geometry.Theory.Line M N) (Ne.{1} Geometry.Theory.Line L N))) (Not (Geometry.Theory.Concurrent L M N)))))\end{verbatim}
\end{theorem}

% === Geometry.Ch2.Prop.P1 ===

\begin{theorem}
  \label{Geometry.Ch2.Prop.P1.direct}
  \uses{Geometry.Theory.Line, Geometry.Theory.Parallel, Geometry.Theory.Point, Geometry.Theory.I1}
  \lean{Geometry.Ch2.Prop.P1.direct}
  \leanok
  pp. 71: If `l` and `m` are distinct lines that are not parallel, then `l` and `m` have a unique point in common 
\begin{verbatim}direct : forall {L : Geometry.Theory.Line} {M : Geometry.Theory.Line}, (Ne.{1} Geometry.Theory.Line L M) -> (Not (Geometry.Theory.Parallel L M)) -> (ExistsUnique.{1} Geometry.Theory.Point (fun (P : Geometry.Theory.Point) => And (Membership.mem.{0, 0} Geometry.Theory.Point Geometry.Theory.Line (Set.instMembership.{0} Geometry.Theory.Point) L P) (Membership.mem.{0, 0} Geometry.Theory.Point Geometry.Theory.Line (Set.instMembership.{0} Geometry.Theory.Point) M P)))\end{verbatim}
\end{theorem}

% === Geometry.Ch3.Ex.Ex1 ===

\begin{theorem}
  \label{Geometry.Ch3.Ex.Ex1.c}
  \uses{Geometry.Ch3.Prop.B4iii}
  \lean{Geometry.Ch3.Ex.Ex1.c}
  \leanok
  (c) Prove the corrolary to B4 Ed. Note that (c) is covered by the B4iii lemma in it's own file. 
\begin{verbatim}c : forall {L : Geometry.Theory.Line} {A : Geometry.Theory.Point} {B : Geometry.Theory.Point} {C : Geometry.Theory.Point}, (And (Not (Membership.mem.{0, 0} Geometry.Theory.Point Geometry.Theory.Line (Set.instMembership.{0} Geometry.Theory.Point) L A)) (And (Not (Membership.mem.{0, 0} Geometry.Theory.Point Geometry.Theory.Line (Set.instMembership.{0} Geometry.Theory.Point) L B)) (Not (Membership.mem.{0, 0} Geometry.Theory.Point Geometry.Theory.Line (Set.instMembership.{0} Geometry.Theory.Point) L C)))) -> (And (Not (Geometry.Theory.SameSide A B L)) (Geometry.Theory.SameSide B C L)) -> (Not (Geometry.Theory.SameSide A C L))\end{verbatim}
\end{theorem}

\begin{theorem}
  \label{Geometry.Ch3.Ex.Ex1.a}
  \uses{Geometry.Theory.Point, Geometry.Theory.Between, Geometry.Theory.Distinct, Geometry.Theory.Betweenness.abc_imp_distinct, Geometry.Theory.Distinct.mk, Geometry.Theory.Distinct.pairwise, Geometry.Theory.Betweenness.absurdity_abc_acb}
  \lean{Geometry.Ch3.Ex.Ex1.a}
  \leanok
  p146. Given A-B-C and A-C-D: (a) Prove that A,B,C, and D are four distinct points (the proof requires an axiom) 
\begin{verbatim}a : forall {A : Geometry.Theory.Point} {B : Geometry.Theory.Point} {C : Geometry.Theory.Point} {D : Geometry.Theory.Point}, (And (Geometry.Theory.Between A B C) (Geometry.Theory.Between A C D)) -> (Geometry.Theory.Distinct.{0} Geometry.Theory.Point (List.cons.{0} Geometry.Theory.Point A (List.cons.{0} Geometry.Theory.Point B (List.cons.{0} Geometry.Theory.Point C (List.cons.{0} Geometry.Theory.Point D (List.nil.{0} Geometry.Theory.Point))))))\end{verbatim}
\end{theorem}

\begin{theorem}
  \label{Geometry.Ch3.Ex.Ex1.b}
  \uses{Geometry.Theory.Point, Geometry.Theory.Between, Geometry.Theory.Collinear, Geometry.Theory.Distinct, Geometry.Ch3.Ex.Ex1.a, Geometry.Theory.Distinct.pairwise, Geometry.Theory.Betweenness.abc_imp_collinear, Geometry.Theory.Line, Geometry.Theory.Collinear.line, Geometry.Theory.Line.equiv, Geometry.Theory.Collinear.mem}
  \lean{Geometry.Ch3.Ex.Ex1.b}
  \leanok
  (b) Prove that A,B,C, and D are collinear 
\begin{verbatim}b : forall {A : Geometry.Theory.Point} {B : Geometry.Theory.Point} {C : Geometry.Theory.Point} {D : Geometry.Theory.Point}, (And (Geometry.Theory.Between A B C) (Geometry.Theory.Between A C D)) -> (Geometry.Theory.Collinear (List.cons.{0} Geometry.Theory.Point A (List.cons.{0} Geometry.Theory.Point B (List.cons.{0} Geometry.Theory.Point C (List.cons.{0} Geometry.Theory.Point D (List.nil.{0} Geometry.Theory.Point))))))\end{verbatim}
\end{theorem}

% === Geometry.Ch3.Prop ===

\begin{theorem}
  \label{Geometry.Ch3.Prop.P2}
  \uses{Geometry.Theory.Point, Geometry.Theory.Line, Geometry.Theory.LineThrough, Geometry.Ch2.Prop.P3, Geometry.Theory.I2, Geometry.Theory.Collinear, Geometry.Theory.Distinct, Geometry.Theory.Between, Geometry.Theory.B2, Geometry.Theory.Distinct.pairwise, Geometry.Theory.Segment, Geometry.Theory.Parallel, Geometry.Theory.Line.seg_has_endpoints.right, Geometry.Theory.B1a, Geometry.Theory.Intersects, Geometry.Theory.Intersection.single_point_of_intersection, Geometry.Theory.Ray, Geometry.Theory.Intersection.lift_seg_ray, Geometry.Theory.Extension, Geometry.Theory.B1b, Geometry.Theory.SameSide, Geometry.Theory.B4ii, Geometry.Theory.Betweenness.splits_commutes, Geometry.Theory.Betweenness.guards_commutes, Geometry.Theory.B4i}
  \lean{Geometry.Ch3.Prop.P2}
  \leanok
  p112. "Every line bounds exactly two half-planes, and these half-planes have no point in common."  B4 is the plane-separation axiom, 3.2 here is on the path toward proving the more useful line-separation property later in 3.4. I've chosen to notate the halfplanes in the theorem as 'Hl' and 'Hr' for 'left' and 'right' half-plane, respectively. 
\begin{verbatim}P2 : forall {A : Geometry.Theory.Point} {B : Geometry.Theory.Point} (L : Geometry.Theory.Line), (Eq.{1} Geometry.Theory.Line L (Geometry.Theory.LineThrough A B)) -> (Ne.{1} Geometry.Theory.Point A B) -> (Exists.{1} (Set.{0} Geometry.Theory.Point) (fun (Hl : Set.{0} Geometry.Theory.Point) => Exists.{1} (Set.{0} Geometry.Theory.Point) (fun (Hr : Set.{0} Geometry.Theory.Point) => And (forall (P : Geometry.Theory.Point), (Membership.mem.{0, 0} Geometry.Theory.Point Geometry.Theory.Line (Set.instMembership.{0} Geometry.Theory.Point) L P) -> (And (Not (Membership.mem.{0, 0} Geometry.Theory.Point (Set.{0} Geometry.Theory.Point) (Set.instMembership.{0} Geometry.Theory.Point) Hl P)) (Not (Membership.mem.{0, 0} Geometry.Theory.Point (Set.{0} Geometry.Theory.Point) (Set.instMembership.{0} Geometry.Theory.Point) Hr P)))) (Eq.{1} (Set.{0} Geometry.Theory.Point) (Inter.inter.{0} (Set.{0} Geometry.Theory.Point) (Set.instInter.{0} Geometry.Theory.Point) Hl Hr) (EmptyCollection.emptyCollection.{0} (Set.{0} Geometry.Theory.Point) (Set.instEmptyCollection.{0} Geometry.Theory.Point))))))\end{verbatim}
\end{theorem}

\begin{theorem}
  \label{Geometry.Ch3.Prop.B4iii}
  \uses{Geometry.Theory.Line, Geometry.Theory.Point, Geometry.Theory.SameSide, Geometry.Theory.B4i, Geometry.Theory.Betweenness.guards_commutes}
  \lean{Geometry.Ch3.Prop.B4iii}
  \leanok
  % TODO: fill in statement
\begin{verbatim}B4iii : forall {L : Geometry.Theory.Line} {A : Geometry.Theory.Point} {B : Geometry.Theory.Point} {C : Geometry.Theory.Point}, (And (Not (Membership.mem.{0, 0} Geometry.Theory.Point Geometry.Theory.Line (Set.instMembership.{0} Geometry.Theory.Point) L A)) (And (Not (Membership.mem.{0, 0} Geometry.Theory.Point Geometry.Theory.Line (Set.instMembership.{0} Geometry.Theory.Point) L B)) (Not (Membership.mem.{0, 0} Geometry.Theory.Point Geometry.Theory.Line (Set.instMembership.{0} Geometry.Theory.Point) L C)))) -> (And (Not (Geometry.Theory.SameSide A B L)) (Geometry.Theory.SameSide B C L)) -> (Not (Geometry.Theory.SameSide A C L))\end{verbatim}
\end{theorem}

% === Geometry.Ch3.Prop.P1 ===

\begin{theorem}
  \label{Geometry.Ch3.Prop.P1.i}
  \uses{Geometry.Theory.Point, Geometry.Theory.Segment, Geometry.Theory.Ray, Geometry.Theory.Line.seg_sub_ray, Geometry.Theory.Line.segment_AB_sub_ray_BA, Geometry.Theory.Between, Geometry.Theory.Distinct, Geometry.Theory.Distinct.mk, Geometry.Theory.Collinear, Geometry.Theory.Line, Geometry.Theory.Line.ray_has_endpoints.left, Geometry.Theory.Line.ray_has_endpoints.right, Geometry.Theory.B3, Geometry.Theory.Extension}
  \lean{Geometry.Ch3.Prop.P1.i}
  \leanok
  p.109, "For any two points A and B: (i) Ray A B ∩ Ray B A = Segment A B ..." 
\begin{verbatim}i : forall {A : Geometry.Theory.Point} {B : Geometry.Theory.Point}, (Ne.{1} Geometry.Theory.Point A B) -> (Eq.{1} (Set.{0} Geometry.Theory.Point) (Geometry.Theory.Segment A B) (Inter.inter.{0} (Set.{0} Geometry.Theory.Point) (Set.instInter.{0} Geometry.Theory.Point) (Geometry.Theory.Ray A B) (Geometry.Theory.Ray B A)))\end{verbatim}
\end{theorem}

\begin{theorem}
  \label{Geometry.Ch3.Prop.P1.ii}
  \uses{Geometry.Theory.Point, Geometry.Theory.Ray, Geometry.Theory.LineThrough, Geometry.Theory.Line.ray_sub_line, Geometry.Theory.Line.commutes, Geometry.Theory.Line.ray_has_endpoints.left, Geometry.Theory.Line.ray_has_endpoints.right, Geometry.Theory.Between, Geometry.Theory.Segment, Geometry.Theory.Line.seg_sub_ray, Geometry.Theory.Extension, Geometry.Theory.B1b}
  \lean{Geometry.Ch3.Prop.P1.ii}
  \leanok
  p.109 "... (ii) Ray A B ∪ Ray B A = LineThrough A B" 
\begin{verbatim}ii : forall {A : Geometry.Theory.Point} {B : Geometry.Theory.Point}, (Ne.{1} Geometry.Theory.Point A B) -> (Eq.{1} (Set.{0} Geometry.Theory.Point) (Union.union.{0} (Set.{0} Geometry.Theory.Point) (Set.instUnion.{0} Geometry.Theory.Point) (Geometry.Theory.Ray A B) (Geometry.Theory.Ray B A)) (Geometry.Theory.LineThrough A B))\end{verbatim}
\end{theorem}

% === Geometry.Ch3.Prop.P3 ===

\begin{theorem}
  \label{Geometry.Ch3.Prop.P3.left}
  \uses{Geometry.Theory.Point, Geometry.Theory.Between, Geometry.Theory.Distinct, Geometry.Ch3.Ex.Ex1.a, Geometry.Theory.Distinct.pairwise, Geometry.Theory.Collinear, Geometry.Ch3.Ex.Ex1.b, Geometry.Theory.Line, Geometry.Theory.Collinear.line, Geometry.Theory.LineThrough, Geometry.Theory.Line.equiv, Geometry.Theory.Collinear.mem, Geometry.Theory.Line.line_has_definition_points.left, Geometry.Theory.Line.line_has_definition_points.right, Geometry.Ch2.Prop.P3, Geometry.Theory.Parallel, Geometry.Theory.Intersection.parallel_intersection_is_empty, Geometry.Theory.Intersects, Geometry.Theory.Intersection.single_point_of_intersection, Geometry.Theory.Intersection.intersection_is_unique, Geometry.Theory.SameSide, Geometry.Theory.Segment, Geometry.Theory.Line.line_trichotomy, Geometry.Theory.Distinct.mk, Geometry.Theory.B3, Geometry.Theory.Line.seg_sub_line, Geometry.Theory.Betweenness.absurdity_abc_bac, Geometry.Theory.Betweenness.absurdity_abc_acb, Geometry.Theory.Intersection.uniq, Geometry.Theory.B4i}
  \lean{Geometry.Ch3.Prop.P3.left}
  \leanok
  p112. Given A - B - C and A - C - D, then B - C - D and A - B - D (see Figure 3.9) 
\begin{verbatim}left : forall (A : Geometry.Theory.Point) (B : Geometry.Theory.Point) (C : Geometry.Theory.Point) (D : Geometry.Theory.Point), (And (Geometry.Theory.Between A B C) (Geometry.Theory.Between A C D)) -> (Geometry.Theory.Between B C D)\end{verbatim}
\end{theorem}

% === Geometry.Theory ===

\begin{lemma}
  \label{Geometry.Theory.I1}
  \lean{Geometry.Theory.I1}
  \leanok
  For any two distinct points P and Q, there exists a unique line L which has P and Q 
\begin{verbatim}I1 : forall (P : Geometry.Theory.Point) (Q : Geometry.Theory.Point), (Ne.{1} Geometry.Theory.Point P Q) -> (ExistsUnique.{1} Geometry.Theory.Line (fun (L : Geometry.Theory.Line) => And (Membership.mem.{0, 0} Geometry.Theory.Point Geometry.Theory.Line (Set.instMembership.{0} Geometry.Theory.Point) L P) (Membership.mem.{0, 0} Geometry.Theory.Point Geometry.Theory.Line (Set.instMembership.{0} Geometry.Theory.Point) L Q)))\end{verbatim}
\end{lemma}

\begin{definition}
  \label{Geometry.Theory.Concurrent}
  \uses{Geometry.Theory.Line, Geometry.Theory.Point}
  \lean{Geometry.Theory.Concurrent}
  \leanok
  p.70 "Three ... lines ... are _concurrent_ if there exists a point incident with all of them"  Ed. Author technically makes this apply to any number of lines, if it ever comes up maybe it's worth a refactor to any finite set of lines? 
\begin{verbatim}Concurrent : Geometry.Theory.Line -> Geometry.Theory.Line -> Geometry.Theory.Line -> Prop\end{verbatim}
\end{definition}

\begin{definition}
  \label{Geometry.Theory.Extension}
  \uses{Geometry.Theory.Point, Geometry.Theory.Between}
  \lean{Geometry.Theory.Extension}
  \leanok
  % TODO: fill in statement
\begin{verbatim}Extension : Geometry.Theory.Point -> Geometry.Theory.Point -> (Set.{0} Geometry.Theory.Point)\end{verbatim}
\end{definition}

\begin{definition}
  \label{Geometry.Theory.Collinear}
  \uses{Geometry.Theory.Point, Geometry.Theory.Line}
  \lean{Geometry.Theory.Collinear}
  \leanok
  % TODO: fill in statement
\begin{verbatim}Collinear : (List.{0} Geometry.Theory.Point) -> Prop\end{verbatim}
\end{definition}

\begin{lemma}
  \label{Geometry.Theory.B4ii}
  \lean{Geometry.Theory.B4ii}
  \leanok
  "... (ii) If A and B are on opposite sides of L and if B and C are opposite sides of L, then A and C are on the same side of L." 
\begin{verbatim}B4ii : forall {A : Geometry.Theory.Point} {B : Geometry.Theory.Point} {C : Geometry.Theory.Point} {L : Geometry.Theory.Line}, (And (Not (Membership.mem.{0, 0} Geometry.Theory.Point Geometry.Theory.Line (Set.instMembership.{0} Geometry.Theory.Point) L A)) (And (Not (Membership.mem.{0, 0} Geometry.Theory.Point Geometry.Theory.Line (Set.instMembership.{0} Geometry.Theory.Point) L B)) (Not (Membership.mem.{0, 0} Geometry.Theory.Point Geometry.Theory.Line (Set.instMembership.{0} Geometry.Theory.Point) L C)))) -> (And (Not (Geometry.Theory.SameSide A B L)) (Not (Geometry.Theory.SameSide B C L))) -> (Geometry.Theory.SameSide A C L)\end{verbatim}
\end{lemma}

\begin{lemma}
  \label{Geometry.Theory.Point}
  \lean{Geometry.Theory.Point}
  \leanok
  A point is the fundamental, opaque type we're working with. 
\begin{verbatim}Point : Type\end{verbatim}
\end{lemma}

\begin{lemma}
  \label{Geometry.Theory.I3}
  \lean{Geometry.Theory.I3}
  \leanok
  There exists three distinct points not on any single line ("There exists three non-collinear points", but without mentioning the undefined notion of collinearity) 
\begin{verbatim}I3 : Exists.{1} Geometry.Theory.Point (fun (A : Geometry.Theory.Point) => Exists.{1} Geometry.Theory.Point (fun (B : Geometry.Theory.Point) => Exists.{1} Geometry.Theory.Point (fun (C : Geometry.Theory.Point) => And (And (Ne.{1} Geometry.Theory.Point A B) (And (Ne.{1} Geometry.Theory.Point A C) (Ne.{1} Geometry.Theory.Point B C))) (forall (L : Geometry.Theory.Line), (Membership.mem.{0, 0} Geometry.Theory.Point Geometry.Theory.Line (Set.instMembership.{0} Geometry.Theory.Point) L A) -> (Membership.mem.{0, 0} Geometry.Theory.Point Geometry.Theory.Line (Set.instMembership.{0} Geometry.Theory.Point) L B) -> (Not (Membership.mem.{0, 0} Geometry.Theory.Point Geometry.Theory.Line (Set.instMembership.{0} Geometry.Theory.Point) L C))))))\end{verbatim}
\end{lemma}

\begin{definition}
  \label{Geometry.Theory.unexpandCollinear}
  \lean{Geometry.Theory.unexpandCollinear}
  \leanok
  % TODO: fill in statement
\begin{verbatim}unexpandCollinear : Lean.PrettyPrinter.Unexpander\end{verbatim}
\end{definition}

\begin{lemma}
  \label{Geometry.Theory.Between}
  \lean{Geometry.Theory.Between}
  \leanok
  % TODO: fill in statement
\begin{verbatim}Between : Geometry.Theory.Point -> Geometry.Theory.Point -> Geometry.Theory.Point -> Prop\end{verbatim}
\end{lemma}

\begin{definition}
  \label{Geometry.Theory.LineThrough}
  \uses{Geometry.Theory.Point, Geometry.Theory.Between}
  \lean{Geometry.Theory.LineThrough}
  \leanok
  % TODO: fill in statement
\begin{verbatim}LineThrough : Geometry.Theory.Point -> Geometry.Theory.Point -> (Set.{0} Geometry.Theory.Point)\end{verbatim}
\end{definition}

\begin{lemma}
  \label{Geometry.Theory.B1a}
  \lean{Geometry.Theory.B1a}
  \leanok
  p.108a "If A - B - C, then A,B,C are distinct points on the same line... 
\begin{verbatim}B1a : forall {A : Geometry.Theory.Point} {B : Geometry.Theory.Point} {C : Geometry.Theory.Point}, (Geometry.Theory.Between A B C) -> (And (Geometry.Theory.Distinct.{0} Geometry.Theory.Point (List.cons.{0} Geometry.Theory.Point A (List.cons.{0} Geometry.Theory.Point B (List.cons.{0} Geometry.Theory.Point C (List.nil.{0} Geometry.Theory.Point))))) (Geometry.Theory.Collinear (List.cons.{0} Geometry.Theory.Point A (List.cons.{0} Geometry.Theory.Point B (List.cons.{0} Geometry.Theory.Point C (List.nil.{0} Geometry.Theory.Point))))))\end{verbatim}
\end{lemma}

\begin{lemma}
  \label{Geometry.Theory.B4i}
  \lean{Geometry.Theory.B4i}
  \leanok
  p.110 "Betweenness Axiom 4 (Plane Separation). For every line L and for any three points A, B, and C not on L: (i) If A and B are on the same side of L and if B and C are on the same side of L, the A and C are on the same side of L..." 
\begin{verbatim}B4i : forall {A : Geometry.Theory.Point} {B : Geometry.Theory.Point} {C : Geometry.Theory.Point} {L : Geometry.Theory.Line}, (And (Not (Membership.mem.{0, 0} Geometry.Theory.Point Geometry.Theory.Line (Set.instMembership.{0} Geometry.Theory.Point) L A)) (And (Not (Membership.mem.{0, 0} Geometry.Theory.Point Geometry.Theory.Line (Set.instMembership.{0} Geometry.Theory.Point) L B)) (Not (Membership.mem.{0, 0} Geometry.Theory.Point Geometry.Theory.Line (Set.instMembership.{0} Geometry.Theory.Point) L C)))) -> (And (Geometry.Theory.SameSide A B L) (Geometry.Theory.SameSide B C L)) -> (Geometry.Theory.SameSide A C L)\end{verbatim}
\end{lemma}

\begin{definition}
  \label{Geometry.Theory.Line}
  \uses{Geometry.Theory.Point}
  \lean{Geometry.Theory.Line}
  \leanok
  Ed: In the text, the author ends up using the 'Line is a Set of points' to define segments, rays, and implicitly uses the intuitive idea ("A line is the set of collinear points that contain at least two known points"). However, Ch2 uses an opaque 'Line' type and reasons only about it's properties without definition. I try to replicate this in my implementation of Ch2, but do define it as a set 'up front' 
\begin{verbatim}Line : Type\end{verbatim}
\end{definition}

\begin{lemma}
  \label{Geometry.Theory.LotEMGuards}
  \lean{Geometry.Theory.LotEMGuards}
  \leanok
  Ed. The author refers to the law of the excluded middle, Lean does not include it by default and generally I want to avoid including it everywhere, this is a limited application of it which should help our purpose. 
\begin{verbatim}LotEMGuards : forall {A : Geometry.Theory.Point} {B : Geometry.Theory.Point} {L : Geometry.Theory.Line}, Or (Not (Geometry.Theory.SameSide A B L)) (Geometry.Theory.SameSide A B L)\end{verbatim}
\end{lemma}

\begin{definition}
  \label{Geometry.Theory.SameSide}
  \uses{Geometry.Theory.Point, Geometry.Theory.Line, Geometry.Theory.Segment}
  \lean{Geometry.Theory.SameSide}
  \leanok
  p.110 "Definition. Let L be any line, and A and B points that do not lie on L. If A = B or if the segment A B contains no points that lie on L, we say that A and B are _on the same side_ of L; whereas, if A ≠ B and segment A B does intersect L, we say that A and B are _on opposite sides_ of L (see Figure 3.6). The law of the excluded middle (Logic Rule 10) tells us that A and B are either on the same side or on opposite sides of L" 
\begin{verbatim}SameSide : Geometry.Theory.Point -> Geometry.Theory.Point -> Geometry.Theory.Line -> Prop\end{verbatim}
\end{definition}

\begin{lemma}
  \label{Geometry.Theory.B3}
  \lean{Geometry.Theory.B3}
  \leanok
  p.108 "If A, B, and C are three distinct points lying on the same line, then one and only one of the points is between the other two." 
\begin{verbatim}B3 : forall (A : Geometry.Theory.Point) (B : Geometry.Theory.Point) (C : Geometry.Theory.Point), (And (Geometry.Theory.Distinct.{0} Geometry.Theory.Point (List.cons.{0} Geometry.Theory.Point A (List.cons.{0} Geometry.Theory.Point B (List.cons.{0} Geometry.Theory.Point C (List.nil.{0} Geometry.Theory.Point))))) (Geometry.Theory.Collinear (List.cons.{0} Geometry.Theory.Point A (List.cons.{0} Geometry.Theory.Point B (List.cons.{0} Geometry.Theory.Point C (List.nil.{0} Geometry.Theory.Point)))))) -> (Or (And (Geometry.Theory.Between A B C) (And (Not (Geometry.Theory.Between B A C)) (Not (Geometry.Theory.Between A C B)))) (Or (And (Not (Geometry.Theory.Between A B C)) (And (Geometry.Theory.Between B A C) (Not (Geometry.Theory.Between A C B)))) (And (Not (Geometry.Theory.Between A B C)) (And (Not (Geometry.Theory.Between B A C)) (Geometry.Theory.Between A C B)))))\end{verbatim}
\end{lemma}

\begin{definition}
  \label{Geometry.Theory.Segment}
  \uses{Geometry.Theory.Point, Geometry.Theory.Between}
  \lean{Geometry.Theory.Segment}
  \leanok
  % TODO: fill in statement
\begin{verbatim}Segment : Geometry.Theory.Point -> Geometry.Theory.Point -> (Set.{0} Geometry.Theory.Point)\end{verbatim}
\end{definition}

\begin{lemma}
  \label{Geometry.Theory.B1b}
  \lean{Geometry.Theory.B1b}
  \leanok
  p.108b ... and [A - B - C iff] C - B - A.""  Ed. Note, I separated these parts of the axiom to make rewriting a bit easier. The author even notes, "The second part (C * B * A) makes the obvious remark that 'betwen A and C' means the same as 'between C and A'" Making it a separate axiom means I won't have to dig it out of the pile of parts that is 1a. 
\begin{verbatim}B1b : forall {A : Geometry.Theory.Point} {B : Geometry.Theory.Point} {C : Geometry.Theory.Point}, Iff (Geometry.Theory.Between A B C) (Geometry.Theory.Between C B A)\end{verbatim}
\end{lemma}

\begin{lemma}
  \label{Geometry.Theory.B2}
  \lean{Geometry.Theory.B2}
  \leanok
  p.108 "Given any two distinct points B and D, there exist points A, C, and E lying on →ₗBD such that A * B * D, B * C * D, and B * D * E".  Ed. I like to call this the 'density' axiom because, used recursively, it posits something like the density of rationals -- for any two distinct points on a line, there is always a point between them. 
\begin{verbatim}B2 : forall (B : Geometry.Theory.Point) (D : Geometry.Theory.Point), (Ne.{1} Geometry.Theory.Point B D) -> (Exists.{1} Geometry.Theory.Point (fun (A : Geometry.Theory.Point) => Exists.{1} Geometry.Theory.Point (fun (C : Geometry.Theory.Point) => Exists.{1} Geometry.Theory.Point (fun (E : Geometry.Theory.Point) => And (Geometry.Theory.Collinear (List.cons.{0} Geometry.Theory.Point A (List.cons.{0} Geometry.Theory.Point B (List.cons.{0} Geometry.Theory.Point C (List.cons.{0} Geometry.Theory.Point D (List.cons.{0} Geometry.Theory.Point E (List.nil.{0} Geometry.Theory.Point))))))) (And (Geometry.Theory.Distinct.{0} Geometry.Theory.Point (List.cons.{0} Geometry.Theory.Point A (List.cons.{0} Geometry.Theory.Point B (List.cons.{0} Geometry.Theory.Point C (List.cons.{0} Geometry.Theory.Point D (List.cons.{0} Geometry.Theory.Point E (List.nil.{0} Geometry.Theory.Point))))))) (And (Geometry.Theory.Between A B D) (And (Geometry.Theory.Between B C D) (Geometry.Theory.Between B D E))))))))\end{verbatim}
\end{lemma}

\begin{definition}
  \label{Geometry.Theory.Parallel}
  \uses{Geometry.Theory.Line, Geometry.Theory.Point}
  \lean{Geometry.Theory.Parallel}
  \leanok
  p. 20, "Two lines `l` and `m` are parallel if they do not intersect, i.e., if no point lies on both of them. We denote this by `l ‖ m`"  p. 70, "Lines `l` and `m` are _parallel_ if they are distinct lines and no point is incident to both of them."  Ed. This gets defined twice, the definitions are equivalent 
\begin{verbatim}Parallel : Geometry.Theory.Line -> Geometry.Theory.Line -> Prop\end{verbatim}
\end{definition}

\begin{definition}
  \label{Geometry.Theory.Intersects}
  \uses{Geometry.Theory.Line, Geometry.Theory.Point}
  \lean{Geometry.Theory.Intersects}
  \leanok
  % TODO: fill in statement
\begin{verbatim}Intersects : Geometry.Theory.Line -> Geometry.Theory.Line -> Geometry.Theory.Point -> Prop\end{verbatim}
\end{definition}

\begin{lemma}
  \label{Geometry.Theory.Distinct}
  \lean{Geometry.Theory.Distinct}
  \leanok
  % TODO: fill in statement
\begin{verbatim}Distinct : forall {α : Type.{u_1}}, (List.{u_1} α) -> Prop\end{verbatim}
\end{lemma}

\begin{lemma}
  \label{Geometry.Theory.I2}
  \lean{Geometry.Theory.I2}
  \leanok
  For any line, there are at least two distinct points on it 
\begin{verbatim}I2 : forall (L : Geometry.Theory.Line), Exists.{1} Geometry.Theory.Point (fun (A : Geometry.Theory.Point) => Exists.{1} Geometry.Theory.Point (fun (B : Geometry.Theory.Point) => And (Ne.{1} Geometry.Theory.Point A B) (And (Membership.mem.{0, 0} Geometry.Theory.Point Geometry.Theory.Line (Set.instMembership.{0} Geometry.Theory.Point) L A) (Membership.mem.{0, 0} Geometry.Theory.Point Geometry.Theory.Line (Set.instMembership.{0} Geometry.Theory.Point) L B))))\end{verbatim}
\end{lemma}

\begin{definition}
  \label{Geometry.Theory.Ray}
  \uses{Geometry.Theory.Point, Geometry.Theory.Segment, Geometry.Theory.Extension}
  \lean{Geometry.Theory.Ray}
  \leanok
  % TODO: fill in statement
\begin{verbatim}Ray : Geometry.Theory.Point -> Geometry.Theory.Point -> (Set.{0} Geometry.Theory.Point)\end{verbatim}
\end{definition}

% === Geometry.Theory.B2 ===

\begin{theorem}
  \label{Geometry.Theory.B2.center}
  \uses{Geometry.Theory.Point, Geometry.Theory.Collinear, Geometry.Theory.Distinct, Geometry.Theory.Between, Geometry.Theory.B2}
  \lean{Geometry.Theory.B2.center}
  \leanok
  Construct a point 'in between' points BD on the induced line B D 
\begin{verbatim}center : forall (B : Geometry.Theory.Point) (D : Geometry.Theory.Point), (Ne.{1} Geometry.Theory.Point B D) -> (Exists.{1} Geometry.Theory.Point (fun (C : Geometry.Theory.Point) => And (Geometry.Theory.Collinear (List.cons.{0} Geometry.Theory.Point B (List.cons.{0} Geometry.Theory.Point C (List.cons.{0} Geometry.Theory.Point D (List.nil.{0} Geometry.Theory.Point))))) (And (Geometry.Theory.Distinct.{0} Geometry.Theory.Point (List.cons.{0} Geometry.Theory.Point B (List.cons.{0} Geometry.Theory.Point C (List.cons.{0} Geometry.Theory.Point D (List.nil.{0} Geometry.Theory.Point))))) (Geometry.Theory.Between B C D))))\end{verbatim}
\end{theorem}

\begin{theorem}
  \label{Geometry.Theory.B2.right}
  \uses{Geometry.Theory.Point, Geometry.Theory.Collinear, Geometry.Theory.Distinct, Geometry.Theory.Between, Geometry.Theory.B2}
  \lean{Geometry.Theory.B2.right}
  \leanok
  Construct a point 'to the right' points BD on the induced line B D 
\begin{verbatim}right : forall (B : Geometry.Theory.Point) (D : Geometry.Theory.Point), (Ne.{1} Geometry.Theory.Point B D) -> (Exists.{1} Geometry.Theory.Point (fun (E : Geometry.Theory.Point) => And (Geometry.Theory.Collinear (List.cons.{0} Geometry.Theory.Point B (List.cons.{0} Geometry.Theory.Point D (List.cons.{0} Geometry.Theory.Point E (List.nil.{0} Geometry.Theory.Point))))) (And (Geometry.Theory.Distinct.{0} Geometry.Theory.Point (List.cons.{0} Geometry.Theory.Point B (List.cons.{0} Geometry.Theory.Point D (List.cons.{0} Geometry.Theory.Point E (List.nil.{0} Geometry.Theory.Point))))) (Geometry.Theory.Between B D E))))\end{verbatim}
\end{theorem}

\begin{theorem}
  \label{Geometry.Theory.B2.left}
  \uses{Geometry.Theory.Point, Geometry.Theory.Collinear, Geometry.Theory.Distinct, Geometry.Theory.Between, Geometry.Theory.B2}
  \lean{Geometry.Theory.B2.left}
  \leanok
  Construct a point 'to the left' of points BD on the induced line B D 
\begin{verbatim}left : forall (B : Geometry.Theory.Point) (D : Geometry.Theory.Point), (Ne.{1} Geometry.Theory.Point B D) -> (Exists.{1} Geometry.Theory.Point (fun (A : Geometry.Theory.Point) => And (Geometry.Theory.Collinear (List.cons.{0} Geometry.Theory.Point A (List.cons.{0} Geometry.Theory.Point B (List.cons.{0} Geometry.Theory.Point D (List.nil.{0} Geometry.Theory.Point))))) (And (Geometry.Theory.Distinct.{0} Geometry.Theory.Point (List.cons.{0} Geometry.Theory.Point A (List.cons.{0} Geometry.Theory.Point B (List.cons.{0} Geometry.Theory.Point D (List.nil.{0} Geometry.Theory.Point))))) (Geometry.Theory.Between A B D))))\end{verbatim}
\end{theorem}

% === Geometry.Theory.Betweenness ===

\begin{theorem}
  \label{Geometry.Theory.Betweenness.abc_imp_collinear}
  \uses{Geometry.Theory.Point, Geometry.Theory.Between, Geometry.Theory.Distinct, Geometry.Theory.Collinear, Geometry.Theory.B1a}
  \lean{Geometry.Theory.Betweenness.abc_imp_collinear}
  \leanok
  betweeness implies collinearity 
\begin{verbatim}abc_imp_collinear : forall {A : Geometry.Theory.Point} {B : Geometry.Theory.Point} {C : Geometry.Theory.Point}, (Geometry.Theory.Between A B C) -> (Geometry.Theory.Collinear (List.cons.{0} Geometry.Theory.Point A (List.cons.{0} Geometry.Theory.Point B (List.cons.{0} Geometry.Theory.Point C (List.nil.{0} Geometry.Theory.Point)))))\end{verbatim}
\end{theorem}

\begin{theorem}
  \label{Geometry.Theory.Betweenness.splits_commutes}
  \uses{Geometry.Theory.Point, Geometry.Theory.Line, Geometry.Theory.SameSide, Geometry.Theory.Segment, Geometry.Theory.Line.segment_AB_eq_segment_BA}
  \lean{Geometry.Theory.Betweenness.splits_commutes}
  \leanok
  a line doesn't care about the order of the points it splits 
\begin{verbatim}splits_commutes : forall {A : Geometry.Theory.Point} {B : Geometry.Theory.Point} {L : Geometry.Theory.Line}, (Not (Geometry.Theory.SameSide A B L)) -> (Not (Geometry.Theory.SameSide B A L))\end{verbatim}
\end{theorem}

\begin{theorem}
  \label{Geometry.Theory.Betweenness.absurdity_abc_bac}
  \uses{Geometry.Theory.Point, Geometry.Theory.Between, Geometry.Theory.Distinct, Geometry.Theory.Collinear, Geometry.Theory.B1a, Geometry.Theory.B3}
  \lean{Geometry.Theory.Betweenness.absurdity_abc_bac}
  \leanok
  With respect to a fixed point, every pair of points can be said to either be 'to the left' or 'to the right' of one another 
\begin{verbatim}absurdity_abc_bac : forall {A : Geometry.Theory.Point} {B : Geometry.Theory.Point} {C : Geometry.Theory.Point}, (And (Geometry.Theory.Between A B C) (Geometry.Theory.Between B A C)) -> False\end{verbatim}
\end{theorem}

\begin{theorem}
  \label{Geometry.Theory.Betweenness.guards_commutes}
  \uses{Geometry.Theory.Point, Geometry.Theory.Line, Geometry.Theory.SameSide, Geometry.Theory.Segment, Geometry.Theory.Line.segment_AB_eq_segment_BA}
  \lean{Geometry.Theory.Betweenness.guards_commutes}
  \leanok
  a line doesn't care about the order of the points it guards 
\begin{verbatim}guards_commutes : forall {A : Geometry.Theory.Point} {B : Geometry.Theory.Point} {L : Geometry.Theory.Line}, (Geometry.Theory.SameSide A B L) -> (Geometry.Theory.SameSide B A L)\end{verbatim}
\end{theorem}

\begin{theorem}
  \label{Geometry.Theory.Betweenness.absurdity_abc_acb}
  \uses{Geometry.Theory.Point, Geometry.Theory.Between, Geometry.Theory.Distinct, Geometry.Theory.Collinear, Geometry.Theory.B1a, Geometry.Theory.B3}
  \lean{Geometry.Theory.Betweenness.absurdity_abc_acb}
  \leanok
  With respect to a fixed point, every pair of points can be said to either be 'to the left' or 'to the right' of one another 
\begin{verbatim}absurdity_abc_acb : forall {A : Geometry.Theory.Point} {B : Geometry.Theory.Point} {C : Geometry.Theory.Point}, (And (Geometry.Theory.Between A B C) (Geometry.Theory.Between A C B)) -> False\end{verbatim}
\end{theorem}

\begin{theorem}
  \label{Geometry.Theory.Betweenness.absurdity_abc_cab}
  \uses{Geometry.Theory.Point, Geometry.Theory.Between, Geometry.Theory.Distinct, Geometry.Theory.Collinear, Geometry.Theory.B1a, Geometry.Theory.B3, Geometry.Theory.B1b}
  \lean{Geometry.Theory.Betweenness.absurdity_abc_cab}
  \leanok
  With respect to a pair of fixed points, another point is either 'to the left' or 'to the right' of the pair 
\begin{verbatim}absurdity_abc_cab : forall {A : Geometry.Theory.Point} {B : Geometry.Theory.Point} {C : Geometry.Theory.Point}, (And (Geometry.Theory.Between A B C) (Geometry.Theory.Between C A B)) -> False\end{verbatim}
\end{theorem}

\begin{theorem}
  \label{Geometry.Theory.Betweenness.abc_imp_distinct}
  \uses{Geometry.Theory.Point, Geometry.Theory.Between, Geometry.Theory.Distinct, Geometry.Theory.Collinear, Geometry.Theory.B1a}
  \lean{Geometry.Theory.Betweenness.abc_imp_distinct}
  \leanok
  betweeness implies distinctness 
\begin{verbatim}abc_imp_distinct : forall {A : Geometry.Theory.Point} {B : Geometry.Theory.Point} {C : Geometry.Theory.Point}, (Geometry.Theory.Between A B C) -> (Geometry.Theory.Distinct.{0} Geometry.Theory.Point (List.cons.{0} Geometry.Theory.Point A (List.cons.{0} Geometry.Theory.Point B (List.cons.{0} Geometry.Theory.Point C (List.nil.{0} Geometry.Theory.Point)))))\end{verbatim}
\end{theorem}

% === Geometry.Theory.Collinear ===

\begin{theorem}
  \label{Geometry.Theory.Collinear.redundancy_irrelevance_ABB}
  \uses{Geometry.Theory.Point, Geometry.Theory.Collinear, Geometry.Theory.Line}
  \lean{Geometry.Theory.Collinear.redundancy_irrelevance_ABB}
  \leanok
  % TODO: fill in statement
\begin{verbatim}redundancy_irrelevance_ABB : forall (A : Geometry.Theory.Point) (B : Geometry.Theory.Point), Iff (Geometry.Theory.Collinear (List.cons.{0} Geometry.Theory.Point A (List.cons.{0} Geometry.Theory.Point B (List.cons.{0} Geometry.Theory.Point B (List.nil.{0} Geometry.Theory.Point))))) (Geometry.Theory.Collinear (List.cons.{0} Geometry.Theory.Point A (List.cons.{0} Geometry.Theory.Point B (List.nil.{0} Geometry.Theory.Point))))\end{verbatim}
\end{theorem}

\begin{theorem}
  \label{Geometry.Theory.Collinear.order_irrelevance}
  \uses{Geometry.Theory.Point, Geometry.Theory.Collinear, Geometry.Theory.Line}
  \lean{Geometry.Theory.Collinear.order_irrelevance}
  \leanok
  Collinearity is independent of list order 
\begin{verbatim}order_irrelevance : forall {S : List.{0} Geometry.Theory.Point} {T : List.{0} Geometry.Theory.Point}, (Geometry.Theory.Collinear S) -> (autoParam.{0} (forall (p : Geometry.Theory.Point), Iff (Membership.mem.{0, 0} Geometry.Theory.Point (List.{0} Geometry.Theory.Point) (List.instMembership.{0} Geometry.Theory.Point) S p) (Membership.mem.{0, 0} Geometry.Theory.Point (List.{0} Geometry.Theory.Point) (List.instMembership.{0} Geometry.Theory.Point) T p)) Geometry.Theory.Collinear.order_irrelevance._auto_1) -> (Geometry.Theory.Collinear T)\end{verbatim}
\end{theorem}

\begin{theorem}
  \label{Geometry.Theory.Collinear.any_two_points_are_collinear}
  \uses{Geometry.Theory.Point, Geometry.Theory.Line, Geometry.Theory.Collinear, Geometry.Theory.I1}
  \lean{Geometry.Theory.Collinear.any_two_points_are_collinear}
  \leanok
  There is a line between any two points, so by definition any two points are collinear 
\begin{verbatim}any_two_points_are_collinear : forall {A : Geometry.Theory.Point} {B : Geometry.Theory.Point}, (Ne.{1} Geometry.Theory.Point A B) -> (Geometry.Theory.Collinear (List.cons.{0} Geometry.Theory.Point A (List.cons.{0} Geometry.Theory.Point B (List.nil.{0} Geometry.Theory.Point))))\end{verbatim}
\end{theorem}

\begin{theorem}
  \label{Geometry.Theory.Collinear.mem}
  \uses{Geometry.Theory.Point, Geometry.Theory.Collinear, Geometry.Theory.Collinear.on_line}
  \lean{Geometry.Theory.Collinear.mem}
  \leanok
  % TODO: fill in statement
\begin{verbatim}mem : forall {points : List.{0} Geometry.Theory.Point} (h : Geometry.Theory.Collinear points) (p : Geometry.Theory.Point), (autoParam.{0} (Membership.mem.{0, 0} Geometry.Theory.Point (List.{0} Geometry.Theory.Point) (List.instMembership.{0} Geometry.Theory.Point) points p) Geometry.Theory.Collinear.mem._auto_1) -> (Membership.mem.{0, 0} Geometry.Theory.Point Geometry.Theory.Line (Set.instMembership.{0} Geometry.Theory.Point) (Geometry.Theory.Collinear.line points h) p)\end{verbatim}
\end{theorem}

\begin{theorem}
  \label{Geometry.Theory.Collinear.on_line}
  \uses{Geometry.Theory.Point, Geometry.Theory.Collinear, Geometry.Theory.Line}
  \lean{Geometry.Theory.Collinear.on_line}
  \leanok
  % TODO: fill in statement
\begin{verbatim}on_line : forall {points : List.{0} Geometry.Theory.Point} (h : Geometry.Theory.Collinear points) (p : Geometry.Theory.Point), (Membership.mem.{0, 0} Geometry.Theory.Point (List.{0} Geometry.Theory.Point) (List.instMembership.{0} Geometry.Theory.Point) points p) -> (Membership.mem.{0, 0} Geometry.Theory.Point Geometry.Theory.Line (Set.instMembership.{0} Geometry.Theory.Point) (Geometry.Theory.Collinear.line points h) p)\end{verbatim}
\end{theorem}

\begin{theorem}
  \label{Geometry.Theory.Collinear.redundancy_irrelevance_BAB}
  \uses{Geometry.Theory.Point, Geometry.Theory.Collinear, Geometry.Theory.Line}
  \lean{Geometry.Theory.Collinear.redundancy_irrelevance_BAB}
  \leanok
  % TODO: fill in statement
\begin{verbatim}redundancy_irrelevance_BAB : forall (A : Geometry.Theory.Point) (B : Geometry.Theory.Point), Iff (Geometry.Theory.Collinear (List.cons.{0} Geometry.Theory.Point B (List.cons.{0} Geometry.Theory.Point A (List.cons.{0} Geometry.Theory.Point B (List.nil.{0} Geometry.Theory.Point))))) (Geometry.Theory.Collinear (List.cons.{0} Geometry.Theory.Point A (List.cons.{0} Geometry.Theory.Point B (List.nil.{0} Geometry.Theory.Point))))\end{verbatim}
\end{theorem}

\begin{definition}
  \label{Geometry.Theory.Collinear.line}
  \uses{Geometry.Theory.Point, Geometry.Theory.Collinear, Geometry.Theory.Line}
  \lean{Geometry.Theory.Collinear.line}
  \leanok
  % TODO: fill in statement
\begin{verbatim}line : forall {points : List.{0} Geometry.Theory.Point}, (Geometry.Theory.Collinear points) -> Geometry.Theory.Line\end{verbatim}
\end{definition}

% === Geometry.Theory.Distinct ===

\begin{theorem}
  \label{Geometry.Theory.Distinct.pairwise}
  \uses{Geometry.Theory.Distinct}
  \lean{Geometry.Theory.Distinct.pairwise}
  \leanok
  % TODO: fill in statement
\begin{verbatim}pairwise : forall {α : Type.{u_1}} {points : List.{u_1} α}, (Geometry.Theory.Distinct.{u_1} α points) -> (List.Pairwise.{u_1} α (fun (x1._@.Geometry.Theory.Distinct.2449265583._hygCtx._hyg.9 : α) (x2._@.Geometry.Theory.Distinct.2449265583._hygCtx._hyg.9 : α) => Ne.{succ u_1} α x1._@.Geometry.Theory.Distinct.2449265583._hygCtx._hyg.9 x2._@.Geometry.Theory.Distinct.2449265583._hygCtx._hyg.9) points)\end{verbatim}
\end{theorem}

\begin{lemma}
  \label{Geometry.Theory.Distinct.mk}
  \lean{Geometry.Theory.Distinct.mk}
  \leanok
  % TODO: fill in statement
\begin{verbatim}mk : forall {α : Type.{u_1}} {points : List.{u_1} α}, (List.Pairwise.{u_1} α (fun (x1._@.Geometry.Theory.Distinct.2449265583._hygCtx._hyg.9 : α) (x2._@.Geometry.Theory.Distinct.2449265583._hygCtx._hyg.9 : α) => Ne.{succ u_1} α x1._@.Geometry.Theory.Distinct.2449265583._hygCtx._hyg.9 x2._@.Geometry.Theory.Distinct.2449265583._hygCtx._hyg.9) points) -> (Geometry.Theory.Distinct.{u_1} α points)\end{verbatim}
\end{lemma}

% === Geometry.Theory.Intersection ===

\begin{theorem}
  \label{Geometry.Theory.Intersection.intersections_are_not_parallel}
  \uses{Geometry.Theory.Line, Geometry.Theory.Point, Geometry.Theory.Intersects, Geometry.Theory.Parallel}
  \lean{Geometry.Theory.Intersection.intersections_are_not_parallel}
  \leanok
  If L intersects M anywhere, then L cannot be parallel to M 
\begin{verbatim}intersections_are_not_parallel : forall {L : Geometry.Theory.Line} {M : Geometry.Theory.Line} {P : Geometry.Theory.Point}, (Geometry.Theory.Intersects L M P) -> (Not (Geometry.Theory.Parallel L M))\end{verbatim}
\end{theorem}

\begin{theorem}
  \label{Geometry.Theory.Intersection.symm}
  \uses{Geometry.Theory.Line, Geometry.Theory.Point, Geometry.Theory.Intersects}
  \lean{Geometry.Theory.Intersection.symm}
  \leanok
  L intersects M is the same as M intersects L. 
\begin{verbatim}symm : forall {L : Geometry.Theory.Line} {M : Geometry.Theory.Line} {X : Geometry.Theory.Point}, Iff (Geometry.Theory.Intersects L M X) (Geometry.Theory.Intersects M L X)\end{verbatim}
\end{theorem}

\begin{theorem}
  \label{Geometry.Theory.Intersection.lift_seg_line}
  \uses{Geometry.Theory.Point, Geometry.Theory.Line, Geometry.Theory.Intersects, Geometry.Theory.Segment, Geometry.Theory.Intersection.lift_ray_line, Geometry.Theory.Intersection.lift_seg_ray}
  \lean{Geometry.Theory.Intersection.lift_seg_line}
  \leanok
  % TODO: fill in statement
\begin{verbatim}lift_seg_line : forall {A : Geometry.Theory.Point} {B : Geometry.Theory.Point} {L : Geometry.Theory.Line} {X : Geometry.Theory.Point} {AneB : Ne.{1} Geometry.Theory.Point A B}, (Geometry.Theory.Intersects L (Geometry.Theory.Segment A B) X) -> (Geometry.Theory.Intersects L (Geometry.Theory.LineThrough A B) X)\end{verbatim}
\end{theorem}

\begin{theorem}
  \label{Geometry.Theory.Intersection.seg_inter_ext_empty}
  \uses{Geometry.Theory.Point, Geometry.Theory.Segment, Geometry.Theory.Extension, Geometry.Theory.Between, Geometry.Theory.Betweenness.absurdity_abc_acb}
  \lean{Geometry.Theory.Intersection.seg_inter_ext_empty}
  \leanok
  No points are contained on the intersection of a segment and it's related extension 
\begin{verbatim}seg_inter_ext_empty : forall {A : Geometry.Theory.Point} {B : Geometry.Theory.Point}, Eq.{1} (Set.{0} Geometry.Theory.Point) (Inter.inter.{0} (Set.{0} Geometry.Theory.Point) (Set.instInter.{0} Geometry.Theory.Point) (Geometry.Theory.Segment A B) (Geometry.Theory.Extension A B)) (EmptyCollection.emptyCollection.{0} (Set.{0} Geometry.Theory.Point) (Set.instEmptyCollection.{0} Geometry.Theory.Point))\end{verbatim}
\end{theorem}

\begin{theorem}
  \label{Geometry.Theory.Intersection.uniq}
  \uses{Geometry.Theory.Line, Geometry.Theory.Point, Geometry.Theory.Intersects}
  \lean{Geometry.Theory.Intersection.uniq}
  \leanok
  If two lines intersect, their intersection is unique. 
\begin{verbatim}uniq : forall {L : Geometry.Theory.Line} {M : Geometry.Theory.Line} {X : Geometry.Theory.Point} {Y : Geometry.Theory.Point}, (And (Geometry.Theory.Intersects L M X) (Geometry.Theory.Intersects L M Y)) -> (Eq.{1} Geometry.Theory.Point X Y)\end{verbatim}
\end{theorem}

\begin{theorem}
  \label{Geometry.Theory.Intersection.inter_touch_left}
  \uses{Geometry.Theory.Line, Geometry.Theory.Point, Geometry.Theory.Intersects}
  \lean{Geometry.Theory.Intersection.inter_touch_left}
  \leanok
  If L intersects M at X, then X is on L 
\begin{verbatim}inter_touch_left : forall {L : Geometry.Theory.Line} {M : Geometry.Theory.Line} {X : Geometry.Theory.Point}, (Geometry.Theory.Intersects L M X) -> (Membership.mem.{0, 0} Geometry.Theory.Point Geometry.Theory.Line (Set.instMembership.{0} Geometry.Theory.Point) L X)\end{verbatim}
\end{theorem}

\begin{theorem}
  \label{Geometry.Theory.Intersection.single_point_of_intersection}
  \uses{Geometry.Theory.Point, Geometry.Theory.Line, Geometry.Theory.Parallel, Geometry.Theory.Intersects, Geometry.Theory.Intersection.intersection_is_unique}
  \lean{Geometry.Theory.Intersection.single_point_of_intersection}
  \leanok
  Intersections of distinct, nonparallel lines contain exactly one point 
\begin{verbatim}single_point_of_intersection : forall (P : Geometry.Theory.Point) (L : Geometry.Theory.Line) (M : Geometry.Theory.Line), (And (Ne.{1} Geometry.Theory.Line L M) (Not (Geometry.Theory.Parallel L M))) -> (Iff (Membership.mem.{0, 0} Geometry.Theory.Point Geometry.Theory.Line (Set.instMembership.{0} Geometry.Theory.Point) (Inter.inter.{0} Geometry.Theory.Line (Set.instInter.{0} Geometry.Theory.Point) L M) P) (Geometry.Theory.Intersects L M P))\end{verbatim}
\end{theorem}

\begin{theorem}
  \label{Geometry.Theory.Intersection.seg_inter_ext_empty'}
  \uses{Geometry.Theory.Point, Geometry.Theory.Segment, Geometry.Theory.Extension, Geometry.Theory.Intersection.seg_inter_ext_empty}
  \lean{Geometry.Theory.Intersection.seg_inter_ext_empty'}
  \leanok
  Points on a segment are not included in the related extension 
\begin{verbatim}seg_inter_ext_empty' : forall {A : Geometry.Theory.Point} {B : Geometry.Theory.Point} {X : Geometry.Theory.Point}, (Membership.mem.{0, 0} Geometry.Theory.Point (Set.{0} Geometry.Theory.Point) (Set.instMembership.{0} Geometry.Theory.Point) (Geometry.Theory.Segment A B) X) -> (Not (Membership.mem.{0, 0} Geometry.Theory.Point (Set.{0} Geometry.Theory.Point) (Set.instMembership.{0} Geometry.Theory.Point) (Geometry.Theory.Extension A B) X))\end{verbatim}
\end{theorem}

\begin{theorem}
  \label{Geometry.Theory.Intersection.intersection_is_unique}
  \uses{Geometry.Theory.Point, Geometry.Theory.Line, Geometry.Theory.Parallel, Geometry.Theory.Intersects, Geometry.Ch2.Prop.P1}
  \lean{Geometry.Theory.Intersection.intersection_is_unique}
  \leanok
  If L and M are distinct, nonparallel lines, and X and Y are found in their intersection, X and Y are equal 
\begin{verbatim}intersection_is_unique : forall {X : Geometry.Theory.Point} {Y : Geometry.Theory.Point} (L : Geometry.Theory.Line) (M : Geometry.Theory.Line), (Ne.{1} Geometry.Theory.Line L M) -> (Not (Geometry.Theory.Parallel L M)) -> (And (Membership.mem.{0, 0} Geometry.Theory.Point Geometry.Theory.Line (Set.instMembership.{0} Geometry.Theory.Point) (Inter.inter.{0} Geometry.Theory.Line (Set.instInter.{0} Geometry.Theory.Point) L M) X) (Membership.mem.{0, 0} Geometry.Theory.Point Geometry.Theory.Line (Set.instMembership.{0} Geometry.Theory.Point) (Inter.inter.{0} Geometry.Theory.Line (Set.instInter.{0} Geometry.Theory.Point) L M) Y)) -> (Eq.{1} Geometry.Theory.Point X Y)\end{verbatim}
\end{theorem}

\begin{theorem}
  \label{Geometry.Theory.Intersection.inter_touch_right}
  \uses{Geometry.Theory.Line, Geometry.Theory.Point, Geometry.Theory.Intersects}
  \lean{Geometry.Theory.Intersection.inter_touch_right}
  \leanok
  If L intersects M at X, then X is on M 
\begin{verbatim}inter_touch_right : forall {L : Geometry.Theory.Line} {M : Geometry.Theory.Line} {X : Geometry.Theory.Point}, (Geometry.Theory.Intersects L M X) -> (Membership.mem.{0, 0} Geometry.Theory.Point Geometry.Theory.Line (Set.instMembership.{0} Geometry.Theory.Point) M X)\end{verbatim}
\end{theorem}

\begin{theorem}
  \label{Geometry.Theory.Intersection.uniq_solitary}
  \uses{Geometry.Theory.Line, Geometry.Theory.Point, Geometry.Theory.Intersects}
  \lean{Geometry.Theory.Intersection.uniq_solitary}
  \leanok
  If L intersects M at X, then forall P not equal to X, if P on L, then P off M. 
\begin{verbatim}uniq_solitary : forall {L : Geometry.Theory.Line} {M : Geometry.Theory.Line} {X : Geometry.Theory.Point}, (And (Ne.{1} Geometry.Theory.Line L M) (Geometry.Theory.Intersects L M X)) -> (forall (P : Geometry.Theory.Point), (And (Ne.{1} Geometry.Theory.Point P X) (Membership.mem.{0, 0} Geometry.Theory.Point Geometry.Theory.Line (Set.instMembership.{0} Geometry.Theory.Point) L P)) -> (Not (Membership.mem.{0, 0} Geometry.Theory.Point Geometry.Theory.Line (Set.instMembership.{0} Geometry.Theory.Point) M P)))\end{verbatim}
\end{theorem}

\begin{theorem}
  \label{Geometry.Theory.Intersection.lift_seg_ray}
  \uses{Geometry.Theory.Point, Geometry.Theory.Line, Geometry.Theory.Intersects, Geometry.Theory.Segment, Geometry.Theory.Intersection.inter_touch_right, Geometry.Theory.Intersection.inter_touch_left, Geometry.Theory.Ray, Geometry.Theory.Extension, Geometry.Theory.Between, Geometry.Theory.Parallel, Geometry.Theory.Intersection.intersection_is_unique}
  \lean{Geometry.Theory.Intersection.lift_seg_ray}
  \leanok
  If a line intersects a segment, then it intersects the ray containing that segment 
\begin{verbatim}lift_seg_ray : forall {A : Geometry.Theory.Point} {B : Geometry.Theory.Point} {L : Geometry.Theory.Line} {X : Geometry.Theory.Point}, (Ne.{1} Geometry.Theory.Point A B) -> (Geometry.Theory.Intersects L (Geometry.Theory.Segment A B) X) -> (Geometry.Theory.Intersects L (Geometry.Theory.Ray A B) X)\end{verbatim}
\end{theorem}

\begin{theorem}
  \label{Geometry.Theory.Intersection.parallel_intersection_is_empty}
  \uses{Geometry.Theory.Line, Geometry.Theory.Parallel, Geometry.Theory.Point}
  \lean{Geometry.Theory.Intersection.parallel_intersection_is_empty}
  \leanok
  If L and M are distinct, parallel lines, their intersection is empty 
\begin{verbatim}parallel_intersection_is_empty : forall (L : Geometry.Theory.Line) (M : Geometry.Theory.Line), (Ne.{1} Geometry.Theory.Line L M) -> (Geometry.Theory.Parallel L M) -> (Eq.{1} Geometry.Theory.Line (Inter.inter.{0} Geometry.Theory.Line (Set.instInter.{0} Geometry.Theory.Point) L M) (EmptyCollection.emptyCollection.{0} Geometry.Theory.Line (Set.instEmptyCollection.{0} Geometry.Theory.Point)))\end{verbatim}
\end{theorem}

\begin{theorem}
  \label{Geometry.Theory.Intersection.lift_ray_line}
  \uses{Geometry.Theory.Point, Geometry.Theory.Line, Geometry.Theory.Intersects, Geometry.Theory.Ray, Geometry.Theory.Intersection.inter_touch_right, Geometry.Theory.Intersection.inter_touch_left, Geometry.Theory.Collinear, Geometry.Theory.Line.all_points_on_a_ray_are_collinear, Geometry.Theory.LineThrough, Geometry.Theory.Line.ray_sub_line, Geometry.Theory.Parallel, Geometry.Theory.Intersection.intersections_are_not_parallel, Geometry.Theory.Line.line_is_bigger_than_ray, Geometry.Theory.Line.line_has_definition_points.left, Geometry.Theory.Line.ray_has_endpoints.left, Geometry.Theory.Line.line_has_definition_points.right, Geometry.Theory.Line.ray_has_endpoints.right, Geometry.Theory.Intersection.single_point_of_intersection, Geometry.Theory.Intersection.intersection_is_unique}
  \lean{Geometry.Theory.Intersection.lift_ray_line}
  \leanok
  If a line intersects a ray, then it intersects the line containing the ray 
\begin{verbatim}lift_ray_line : forall {A : Geometry.Theory.Point} {B : Geometry.Theory.Point} {L : Geometry.Theory.Line} {X : Geometry.Theory.Point} {AneB : Ne.{1} Geometry.Theory.Point A B}, (Geometry.Theory.Intersects L (Geometry.Theory.Ray A B) X) -> (Geometry.Theory.Intersects L (Geometry.Theory.LineThrough A B) X)\end{verbatim}
\end{theorem}

\begin{theorem}
  \label{Geometry.Theory.Intersection.ext_inter_seg_empty}
  \uses{Geometry.Theory.Point, Geometry.Theory.Extension, Geometry.Theory.Segment, Geometry.Theory.Intersection.seg_inter_ext_empty}
  \lean{Geometry.Theory.Intersection.ext_inter_seg_empty}
  \leanok
  Points on an extension are off the related segment 
\begin{verbatim}ext_inter_seg_empty : forall {A : Geometry.Theory.Point} {B : Geometry.Theory.Point} {X : Geometry.Theory.Point}, (Membership.mem.{0, 0} Geometry.Theory.Point (Set.{0} Geometry.Theory.Point) (Set.instMembership.{0} Geometry.Theory.Point) (Geometry.Theory.Extension A B) X) -> (Not (Membership.mem.{0, 0} Geometry.Theory.Point (Set.{0} Geometry.Theory.Point) (Set.instMembership.{0} Geometry.Theory.Point) (Geometry.Theory.Segment A B) X))\end{verbatim}
\end{theorem}

% === Geometry.Theory.Line ===

\begin{theorem}
  \label{Geometry.Theory.Line.seg_sub_ray}
  \uses{Geometry.Theory.Point, Geometry.Theory.Segment, Geometry.Theory.Extension}
  \lean{Geometry.Theory.Line.seg_sub_ray}
  \leanok
  pXX "By the definition of segment and ray, `the segment A B ⊆ the ray A B`" 
\begin{verbatim}seg_sub_ray : forall {A : Geometry.Theory.Point} {B : Geometry.Theory.Point}, HasSubset.Subset.{0} (Set.{0} Geometry.Theory.Point) (Set.instHasSubset.{0} Geometry.Theory.Point) (Geometry.Theory.Segment A B) (Geometry.Theory.Ray A B)\end{verbatim}
\end{theorem}

\begin{theorem}
  \label{Geometry.Theory.Line.intersecting_lines_are_not_parallel}
  \uses{Geometry.Theory.Line, Geometry.Theory.Point, Geometry.Theory.Parallel}
  \lean{Geometry.Theory.Line.intersecting_lines_are_not_parallel}
  \leanok
  We need to be able to establish that two intersecting lines are never parallel 
\begin{verbatim}intersecting_lines_are_not_parallel : forall {L : Geometry.Theory.Line} {M : Geometry.Theory.Line} {P : Geometry.Theory.Point}, (Membership.mem.{0, 0} Geometry.Theory.Point Geometry.Theory.Line (Set.instMembership.{0} Geometry.Theory.Point) L P) -> (Membership.mem.{0, 0} Geometry.Theory.Point Geometry.Theory.Line (Set.instMembership.{0} Geometry.Theory.Point) M P) -> (Not (Geometry.Theory.Parallel L M))\end{verbatim}
\end{theorem}

\begin{theorem}
  \label{Geometry.Theory.Line.ABP_imp_P_on_ext_AB}
  \uses{Geometry.Theory.Point, Geometry.Theory.Between, Geometry.Theory.Extension, Geometry.Theory.B1b}
  \lean{Geometry.Theory.Line.ABP_imp_P_on_ext_AB}
  \leanok
  % TODO: fill in statement
\begin{verbatim}ABP_imp_P_on_ext_AB : forall {P : Geometry.Theory.Point} {A : Geometry.Theory.Point} {B : Geometry.Theory.Point}, (And (Ne.{1} Geometry.Theory.Point P A) (Ne.{1} Geometry.Theory.Point P B)) -> (Geometry.Theory.Between A B P) -> (Membership.mem.{0, 0} Geometry.Theory.Point (Set.{0} Geometry.Theory.Point) (Set.instMembership.{0} Geometry.Theory.Point) (Geometry.Theory.Extension A B) P)\end{verbatim}
\end{theorem}

\begin{theorem}
  \label{Geometry.Theory.Line.line_trichotomy}
  \uses{Geometry.Theory.Point, Geometry.Theory.Parallel, Geometry.Theory.Line, Geometry.Ch2.Prop.P1}
  \lean{Geometry.Theory.Line.line_trichotomy}
  \leanok
  An intersection is either empty, a singleton, or the lines are equal. 
\begin{verbatim}line_trichotomy : forall (L : Set.{0} Geometry.Theory.Point) (M : Set.{0} Geometry.Theory.Point), Or (Eq.{1} (Set.{0} Geometry.Theory.Point) (Inter.inter.{0} (Set.{0} Geometry.Theory.Point) (Set.instInter.{0} Geometry.Theory.Point) L M) (EmptyCollection.emptyCollection.{0} (Set.{0} Geometry.Theory.Point) (Set.instEmptyCollection.{0} Geometry.Theory.Point))) (Or (ExistsUnique.{1} Geometry.Theory.Point (fun (X : Geometry.Theory.Point) => Eq.{1} (Set.{0} Geometry.Theory.Point) (Inter.inter.{0} (Set.{0} Geometry.Theory.Point) (Set.instInter.{0} Geometry.Theory.Point) L M) (Singleton.singleton.{0, 0} Geometry.Theory.Point (Set.{0} Geometry.Theory.Point) (Set.instSingletonSet.{0} Geometry.Theory.Point) X))) (Eq.{1} (Set.{0} Geometry.Theory.Point) L M))\end{verbatim}
\end{theorem}

\begin{theorem}
  \label{Geometry.Theory.Line.all_points_on_an_extension_are_collinear}
  \uses{Geometry.Theory.Point, Geometry.Theory.Extension, Geometry.Theory.Betweenness.abc_imp_collinear, Geometry.Theory.Between}
  \lean{Geometry.Theory.Line.all_points_on_an_extension_are_collinear}
  \leanok
  All points on a extension are collinear 
\begin{verbatim}all_points_on_an_extension_are_collinear : forall {P : Geometry.Theory.Point} {A : Geometry.Theory.Point} {B : Geometry.Theory.Point}, (Membership.mem.{0, 0} Geometry.Theory.Point (Set.{0} Geometry.Theory.Point) (Set.instMembership.{0} Geometry.Theory.Point) (Geometry.Theory.Extension A B) P) -> (Geometry.Theory.Collinear (List.cons.{0} Geometry.Theory.Point A (List.cons.{0} Geometry.Theory.Point B (List.cons.{0} Geometry.Theory.Point P (List.nil.{0} Geometry.Theory.Point)))))\end{verbatim}
\end{theorem}

\begin{theorem}
  \label{Geometry.Theory.Line.APB_imp_P_on_ray_AB}
  \uses{Geometry.Theory.Point, Geometry.Theory.Between, Geometry.Theory.Ray, Geometry.Theory.Segment, Geometry.Theory.Extension}
  \lean{Geometry.Theory.Line.APB_imp_P_on_ray_AB}
  \leanok
  % TODO: fill in statement
\begin{verbatim}APB_imp_P_on_ray_AB : forall {P : Geometry.Theory.Point} {A : Geometry.Theory.Point} {B : Geometry.Theory.Point}, (And (Ne.{1} Geometry.Theory.Point P A) (Ne.{1} Geometry.Theory.Point P B)) -> (Geometry.Theory.Between A P B) -> (Membership.mem.{0, 0} Geometry.Theory.Point (Set.{0} Geometry.Theory.Point) (Set.instMembership.{0} Geometry.Theory.Point) (Geometry.Theory.Ray A B) P)\end{verbatim}
\end{theorem}

\begin{theorem}
  \label{Geometry.Theory.Line.commutes}
  \uses{Geometry.Theory.Point, Geometry.Theory.LineThrough, Geometry.Theory.Between}
  \lean{Geometry.Theory.Line.commutes}
  \leanok
  % TODO: fill in statement
\begin{verbatim}commutes : forall {A : Geometry.Theory.Point} {B : Geometry.Theory.Point} {AneB : Ne.{1} Geometry.Theory.Point A B}, Eq.{1} (Set.{0} Geometry.Theory.Point) (Geometry.Theory.LineThrough A B) (Geometry.Theory.LineThrough B A)\end{verbatim}
\end{theorem}

\begin{theorem}
  \label{Geometry.Theory.Line.line_has_definition_points}
  \uses{Geometry.Theory.Point, Geometry.Theory.LineThrough, Geometry.Theory.Line.line_has_definition_points.left, Geometry.Theory.Line.line_has_definition_points.right}
  \lean{Geometry.Theory.Line.line_has_definition_points}
  \leanok
  A line contains the points that define it 
\begin{verbatim}line_has_definition_points : forall {A : Geometry.Theory.Point} {B : Geometry.Theory.Point}, And (Membership.mem.{0, 0} Geometry.Theory.Point (Set.{0} Geometry.Theory.Point) (Set.instMembership.{0} Geometry.Theory.Point) (Geometry.Theory.LineThrough A B) A) (Membership.mem.{0, 0} Geometry.Theory.Point (Set.{0} Geometry.Theory.Point) (Set.instMembership.{0} Geometry.Theory.Point) (Geometry.Theory.LineThrough A B) B)\end{verbatim}
\end{theorem}

\begin{theorem}
  \label{Geometry.Theory.Line.ray_sub_line}
  \uses{Geometry.Theory.Point, Geometry.Theory.Ray, Geometry.Theory.LineThrough, Geometry.Theory.Between, Geometry.Theory.Segment, Geometry.Theory.Extension}
  \lean{Geometry.Theory.Line.ray_sub_line}
  \leanok
  A ray A B is a subset of the line A B 
\begin{verbatim}ray_sub_line : forall {A : Geometry.Theory.Point} {B : Geometry.Theory.Point}, HasSubset.Subset.{0} (Set.{0} Geometry.Theory.Point) (Set.instHasSubset.{0} Geometry.Theory.Point) (Geometry.Theory.Ray A B) (Geometry.Theory.LineThrough A B)\end{verbatim}
\end{theorem}

\begin{theorem}
  \label{Geometry.Theory.Line.APB_imp_P_on_segment_AB}
  \uses{Geometry.Theory.Point, Geometry.Theory.Between, Geometry.Theory.Segment}
  \lean{Geometry.Theory.Line.APB_imp_P_on_segment_AB}
  \leanok
  % TODO: fill in statement
\begin{verbatim}APB_imp_P_on_segment_AB : forall {P : Geometry.Theory.Point} {A : Geometry.Theory.Point} {B : Geometry.Theory.Point}, (And (Ne.{1} Geometry.Theory.Point P A) (Ne.{1} Geometry.Theory.Point P B)) -> (Geometry.Theory.Between A P B) -> (Membership.mem.{0, 0} Geometry.Theory.Point (Set.{0} Geometry.Theory.Point) (Set.instMembership.{0} Geometry.Theory.Point) (Geometry.Theory.Segment A B) P)\end{verbatim}
\end{theorem}

\begin{theorem}
  \label{Geometry.Theory.Line.line_by_definition}
  \uses{Geometry.Theory.Line, Geometry.Theory.Point}
  \lean{Geometry.Theory.Line.line_by_definition}
  \leanok
  A line is the set of all points on it 
\begin{verbatim}line_by_definition : forall (L : Geometry.Theory.Line), Eq.{1} Geometry.Theory.Line L (setOf.{0} Geometry.Theory.Point (fun (P : Geometry.Theory.Point) => Membership.mem.{0, 0} Geometry.Theory.Point Geometry.Theory.Line (Set.instMembership.{0} Geometry.Theory.Point) L P))\end{verbatim}
\end{theorem}

\begin{theorem}
  \label{Geometry.Theory.Line.segment_AB_sub_ray_BA}
  \uses{Geometry.Theory.Point, Geometry.Theory.Segment, Geometry.Theory.Extension, Geometry.Theory.Between, Geometry.Theory.Line.segment_AB_eq_segment_BA}
  \lean{Geometry.Theory.Line.segment_AB_sub_ray_BA}
  \leanok
  The endpoint B is in common here. 
\begin{verbatim}segment_AB_sub_ray_BA : forall {A : Geometry.Theory.Point} {B : Geometry.Theory.Point}, HasSubset.Subset.{0} (Set.{0} Geometry.Theory.Point) (Set.instHasSubset.{0} Geometry.Theory.Point) (Geometry.Theory.Segment A B) (Geometry.Theory.Ray B A)\end{verbatim}
\end{theorem}

\begin{theorem}
  \label{Geometry.Theory.Line.distinguishing_point}
  \uses{Geometry.Theory.Line, Geometry.Theory.Point, Geometry.Theory.Line.coincidence_is_coincidence_of_all_points}
  \lean{Geometry.Theory.Line.distinguishing_point}
  \leanok
  Two lines are distinct iff they have at least one point not in common 
\begin{verbatim}distinguishing_point : forall (L : Geometry.Theory.Line) (M : Geometry.Theory.Line), Iff (Ne.{1} Geometry.Theory.Line L M) (Exists.{1} Geometry.Theory.Point (fun (P : Geometry.Theory.Point) => Or (And (Membership.mem.{0, 0} Geometry.Theory.Point Geometry.Theory.Line (Set.instMembership.{0} Geometry.Theory.Point) L P) (Not (Membership.mem.{0, 0} Geometry.Theory.Point Geometry.Theory.Line (Set.instMembership.{0} Geometry.Theory.Point) M P))) (And (Not (Membership.mem.{0, 0} Geometry.Theory.Point Geometry.Theory.Line (Set.instMembership.{0} Geometry.Theory.Point) L P)) (Membership.mem.{0, 0} Geometry.Theory.Point Geometry.Theory.Line (Set.instMembership.{0} Geometry.Theory.Point) M P))))\end{verbatim}
\end{theorem}

\begin{theorem}
  \label{Geometry.Theory.Line.seg_sub_line}
  \uses{Geometry.Theory.Point, Geometry.Theory.Segment, Geometry.Theory.Ray, Geometry.Theory.Line.seg_sub_ray, Geometry.Theory.LineThrough, Geometry.Theory.Line.ray_sub_line, Geometry.Theory.Between}
  \lean{Geometry.Theory.Line.seg_sub_line}
  \leanok
  A segment is a subset of the line A B 
\begin{verbatim}seg_sub_line : forall {A : Geometry.Theory.Point} {B : Geometry.Theory.Point}, HasSubset.Subset.{0} (Set.{0} Geometry.Theory.Point) (Set.instHasSubset.{0} Geometry.Theory.Point) (Geometry.Theory.Segment A B) (Geometry.Theory.LineThrough A B)\end{verbatim}
\end{theorem}

\begin{theorem}
  \label{Geometry.Theory.Line.segment_AB_eq_segment_BA}
  \uses{Geometry.Theory.Point, Geometry.Theory.Segment, Geometry.Theory.Between, Geometry.Theory.B1b}
  \lean{Geometry.Theory.Line.segment_AB_eq_segment_BA}
  \leanok
  % TODO: fill in statement
\begin{verbatim}segment_AB_eq_segment_BA : forall {A : Geometry.Theory.Point} {B : Geometry.Theory.Point}, Eq.{1} (Set.{0} Geometry.Theory.Point) (Geometry.Theory.Segment A B) (Geometry.Theory.Segment B A)\end{verbatim}
\end{theorem}

\begin{theorem}
  \label{Geometry.Theory.Line.all_points_on_a_segment_are_collinear}
  \uses{Geometry.Theory.Point, Geometry.Theory.Segment, Geometry.Theory.Line.all_points_on_a_line_are_collinear, Geometry.Theory.Line.seg_sub_line}
  \lean{Geometry.Theory.Line.all_points_on_a_segment_are_collinear}
  \leanok
  All points on a segment are collinear 
\begin{verbatim}all_points_on_a_segment_are_collinear : forall {A : Geometry.Theory.Point} {B : Geometry.Theory.Point} {P : Geometry.Theory.Point} {AneB : Ne.{1} Geometry.Theory.Point A B}, (Membership.mem.{0, 0} Geometry.Theory.Point (Set.{0} Geometry.Theory.Point) (Set.instMembership.{0} Geometry.Theory.Point) (Geometry.Theory.Segment A B) P) -> (Geometry.Theory.Collinear (List.cons.{0} Geometry.Theory.Point A (List.cons.{0} Geometry.Theory.Point B (List.cons.{0} Geometry.Theory.Point P (List.nil.{0} Geometry.Theory.Point)))))\end{verbatim}
\end{theorem}

\begin{theorem}
  \label{Geometry.Theory.Line.all_points_on_a_ray_are_collinear}
  \uses{Geometry.Theory.Point, Geometry.Theory.Ray, Geometry.Theory.Line.all_points_on_a_line_are_collinear, Geometry.Theory.Line.ray_sub_line}
  \lean{Geometry.Theory.Line.all_points_on_a_ray_are_collinear}
  \leanok
  All points on a ray are collinear 
\begin{verbatim}all_points_on_a_ray_are_collinear : forall {A : Geometry.Theory.Point} {B : Geometry.Theory.Point} {P : Geometry.Theory.Point} {AneB : Ne.{1} Geometry.Theory.Point A B}, (Membership.mem.{0, 0} Geometry.Theory.Point (Set.{0} Geometry.Theory.Point) (Set.instMembership.{0} Geometry.Theory.Point) (Geometry.Theory.Ray A B) P) -> (Geometry.Theory.Collinear (List.cons.{0} Geometry.Theory.Point A (List.cons.{0} Geometry.Theory.Point B (List.cons.{0} Geometry.Theory.Point P (List.nil.{0} Geometry.Theory.Point)))))\end{verbatim}
\end{theorem}

\begin{theorem}
  \label{Geometry.Theory.Line.concurrence_of_three_lines_is_unique}
  \uses{Geometry.Theory.Line, Geometry.Theory.Concurrent, Geometry.Theory.Point, Geometry.Theory.I1}
  \lean{Geometry.Theory.Line.concurrence_of_three_lines_is_unique}
  \leanok
  Author suggests a lemma, "... to prove it, I could first prove a lemma that if three lines are concurrent, the point at which they meet is unique." p.71 
\begin{verbatim}concurrence_of_three_lines_is_unique : forall {L : Geometry.Theory.Line} {M : Geometry.Theory.Line} {N : Geometry.Theory.Line}, (And (Ne.{1} Geometry.Theory.Line L M) (And (Ne.{1} Geometry.Theory.Line M N) (Ne.{1} Geometry.Theory.Line L N))) -> (Geometry.Theory.Concurrent L M N) -> (ExistsUnique.{1} Geometry.Theory.Point (fun (P : Geometry.Theory.Point) => And (Membership.mem.{0, 0} Geometry.Theory.Point Geometry.Theory.Line (Set.instMembership.{0} Geometry.Theory.Point) L P) (And (Membership.mem.{0, 0} Geometry.Theory.Point Geometry.Theory.Line (Set.instMembership.{0} Geometry.Theory.Point) M P) (Membership.mem.{0, 0} Geometry.Theory.Point Geometry.Theory.Line (Set.instMembership.{0} Geometry.Theory.Point) N P))))\end{verbatim}
\end{theorem}

\begin{theorem}
  \label{Geometry.Theory.Line.equiv}
  \uses{Geometry.Theory.Line, Geometry.Theory.Point, Geometry.Theory.Line.line_trichotomy}
  \lean{Geometry.Theory.Line.equiv}
  \leanok
  If two distinct points are found on two lines, those lines are equal. 
\begin{verbatim}equiv : forall {L : Geometry.Theory.Line} {M : Geometry.Theory.Line} {A : Geometry.Theory.Point} {B : Geometry.Theory.Point}, (Ne.{1} Geometry.Theory.Point A B) -> (And (Membership.mem.{0, 0} Geometry.Theory.Point Geometry.Theory.Line (Set.instMembership.{0} Geometry.Theory.Point) L A) (And (Membership.mem.{0, 0} Geometry.Theory.Point Geometry.Theory.Line (Set.instMembership.{0} Geometry.Theory.Point) M A) (And (Membership.mem.{0, 0} Geometry.Theory.Point Geometry.Theory.Line (Set.instMembership.{0} Geometry.Theory.Point) L B) (Membership.mem.{0, 0} Geometry.Theory.Point Geometry.Theory.Line (Set.instMembership.{0} Geometry.Theory.Point) M B)))) -> (Eq.{1} Geometry.Theory.Line L M)\end{verbatim}
\end{theorem}

\begin{theorem}
  \label{Geometry.Theory.Line.segment_int_extension_is_empty}
  \uses{Geometry.Theory.Point, Geometry.Theory.Segment, Geometry.Theory.Extension, Geometry.Theory.Between, Geometry.Theory.Betweenness.absurdity_abc_acb}
  \lean{Geometry.Theory.Line.segment_int_extension_is_empty}
  \leanok
  % TODO: fill in statement
\begin{verbatim}segment_int_extension_is_empty : forall {A : Geometry.Theory.Point} {B : Geometry.Theory.Point}, Eq.{1} (Set.{0} Geometry.Theory.Point) (Inter.inter.{0} (Set.{0} Geometry.Theory.Point) (Set.instInter.{0} Geometry.Theory.Point) (Geometry.Theory.Segment A B) (Geometry.Theory.Extension A B)) (EmptyCollection.emptyCollection.{0} (Set.{0} Geometry.Theory.Point) (Set.instEmptyCollection.{0} Geometry.Theory.Point))\end{verbatim}
\end{theorem}

\begin{theorem}
  \label{Geometry.Theory.Line.line_is_bigger_than_ray}
  \uses{Geometry.Theory.Line, Geometry.Theory.Point, Geometry.Theory.Ray, Geometry.Theory.Collinear, Geometry.Theory.Distinct, Geometry.Theory.Between, Geometry.Theory.B2.left, Geometry.Theory.Distinct.pairwise, Geometry.Theory.Collinear.line, Geometry.Theory.Line.equiv, Geometry.Theory.Line.ray_has_endpoints.left, Geometry.Theory.Collinear.mem, Geometry.Theory.Line.ray_has_endpoints.right, Geometry.Theory.Segment, Geometry.Theory.Extension, Geometry.Theory.Betweenness.absurdity_abc_bac, Geometry.Theory.Betweenness.absurdity_abc_cab}
  \lean{Geometry.Theory.Line.line_is_bigger_than_ray}
  \leanok
  A line is 'bigger' than a ray in the same way that a line is bigger than a segment 
\begin{verbatim}line_is_bigger_than_ray : forall (L : Geometry.Theory.Line) (A : Geometry.Theory.Point) (B : Geometry.Theory.Point), (Ne.{1} Geometry.Theory.Point A B) -> (Ne.{1} (Set.{0} Geometry.Theory.Point) (Geometry.Theory.Ray A B) L)\end{verbatim}
\end{theorem}

\begin{theorem}
  \label{Geometry.Theory.Line.coincidence_is_coincidence_of_all_points}
  \uses{Geometry.Theory.Line, Geometry.Theory.Point, Geometry.Theory.I2, Geometry.Theory.I1}
  \lean{Geometry.Theory.Line.coincidence_is_coincidence_of_all_points}
  \leanok
  Two lines are coincident iff every point on one is on the other. 
\begin{verbatim}coincidence_is_coincidence_of_all_points : forall (L : Geometry.Theory.Line) (M : Geometry.Theory.Line), Iff (Eq.{1} Geometry.Theory.Line L M) (forall (P : Geometry.Theory.Point), Iff (Membership.mem.{0, 0} Geometry.Theory.Point Geometry.Theory.Line (Set.instMembership.{0} Geometry.Theory.Point) L P) (Membership.mem.{0, 0} Geometry.Theory.Point Geometry.Theory.Line (Set.instMembership.{0} Geometry.Theory.Point) M P))\end{verbatim}
\end{theorem}

\begin{theorem}
  \label{Geometry.Theory.Line.all_points_on_a_line_are_collinear}
  \uses{Geometry.Theory.Point, Geometry.Theory.LineThrough, Geometry.Theory.Between, Geometry.Theory.Collinear, Geometry.Theory.Collinear.redundancy_irrelevance_BAB, Geometry.Theory.Collinear.any_two_points_are_collinear, Geometry.Theory.Collinear.redundancy_irrelevance_ABB, Geometry.Theory.Collinear.order_irrelevance, Geometry.Theory.Betweenness.abc_imp_collinear}
  \lean{Geometry.Theory.Line.all_points_on_a_line_are_collinear}
  \leanok
  All points on a line are collinear 
\begin{verbatim}all_points_on_a_line_are_collinear : forall {A : Geometry.Theory.Point} {B : Geometry.Theory.Point} {P : Geometry.Theory.Point} {AneB : Ne.{1} Geometry.Theory.Point A B}, (Membership.mem.{0, 0} Geometry.Theory.Point (Set.{0} Geometry.Theory.Point) (Set.instMembership.{0} Geometry.Theory.Point) (Geometry.Theory.LineThrough A B) P) -> (Geometry.Theory.Collinear (List.cons.{0} Geometry.Theory.Point A (List.cons.{0} Geometry.Theory.Point B (List.cons.{0} Geometry.Theory.Point P (List.nil.{0} Geometry.Theory.Point)))))\end{verbatim}
\end{theorem}

% === Geometry.Theory.Line.line_has_definition_points ===

\begin{theorem}
  \label{Geometry.Theory.Line.line_has_definition_points.right}
  \uses{Geometry.Theory.Point, Geometry.Theory.Line.ray_sub_line, Geometry.Theory.Line.ray_has_endpoints.right}
  \lean{Geometry.Theory.Line.line_has_definition_points.right}
  \leanok
  A line contains the points that define it 
\begin{verbatim}right : forall {A : Geometry.Theory.Point} {B : Geometry.Theory.Point}, Membership.mem.{0, 0} Geometry.Theory.Point (Set.{0} Geometry.Theory.Point) (Set.instMembership.{0} Geometry.Theory.Point) (Geometry.Theory.LineThrough A B) B\end{verbatim}
\end{theorem}

\begin{theorem}
  \label{Geometry.Theory.Line.line_has_definition_points.left}
  \uses{Geometry.Theory.Point, Geometry.Theory.Line.ray_sub_line, Geometry.Theory.Line.ray_has_endpoints.left}
  \lean{Geometry.Theory.Line.line_has_definition_points.left}
  \leanok
  A line contains the points that define it 
\begin{verbatim}left : forall {A : Geometry.Theory.Point} {B : Geometry.Theory.Point}, Membership.mem.{0, 0} Geometry.Theory.Point (Set.{0} Geometry.Theory.Point) (Set.instMembership.{0} Geometry.Theory.Point) (Geometry.Theory.LineThrough A B) A\end{verbatim}
\end{theorem}

% === Geometry.Theory.Line.ray_has_endpoints ===

\begin{theorem}
  \label{Geometry.Theory.Line.ray_has_endpoints.right}
  \uses{Geometry.Theory.Point, Geometry.Theory.Segment, Geometry.Theory.Extension, Geometry.Theory.Between}
  \lean{Geometry.Theory.Line.ray_has_endpoints.right}
  \leanok
  A ray contains the points that define it 
\begin{verbatim}right : forall {A : Geometry.Theory.Point} {B : Geometry.Theory.Point}, Membership.mem.{0, 0} Geometry.Theory.Point (Set.{0} Geometry.Theory.Point) (Set.instMembership.{0} Geometry.Theory.Point) (Geometry.Theory.Ray A B) B\end{verbatim}
\end{theorem}

\begin{theorem}
  \label{Geometry.Theory.Line.ray_has_endpoints.left}
  \uses{Geometry.Theory.Point, Geometry.Theory.Segment, Geometry.Theory.Extension, Geometry.Theory.Between}
  \lean{Geometry.Theory.Line.ray_has_endpoints.left}
  \leanok
  A ray contains the points that define it 
\begin{verbatim}left : forall {A : Geometry.Theory.Point} {B : Geometry.Theory.Point}, Membership.mem.{0, 0} Geometry.Theory.Point (Set.{0} Geometry.Theory.Point) (Set.instMembership.{0} Geometry.Theory.Point) (Geometry.Theory.Ray A B) A\end{verbatim}
\end{theorem}

% === Geometry.Theory.Line.seg_has_endpoints ===

\begin{theorem}
  \label{Geometry.Theory.Line.seg_has_endpoints.right}
  \uses{Geometry.Theory.Point, Geometry.Theory.Between}
  \lean{Geometry.Theory.Line.seg_has_endpoints.right}
  \leanok
  A segment contains the points that define it 
\begin{verbatim}right : forall {A : Geometry.Theory.Point} {B : Geometry.Theory.Point}, Membership.mem.{0, 0} Geometry.Theory.Point (Set.{0} Geometry.Theory.Point) (Set.instMembership.{0} Geometry.Theory.Point) (Geometry.Theory.Segment A B) B\end{verbatim}
\end{theorem}

\begin{theorem}
  \label{Geometry.Theory.Line.seg_has_endpoints.left}
  \uses{Geometry.Theory.Point, Geometry.Theory.Between}
  \lean{Geometry.Theory.Line.seg_has_endpoints.left}
  \leanok
  A segment contains the points that define it 
\begin{verbatim}left : forall {A : Geometry.Theory.Point} {B : Geometry.Theory.Point}, Membership.mem.{0, 0} Geometry.Theory.Point (Set.{0} Geometry.Theory.Point) (Set.instMembership.{0} Geometry.Theory.Point) (Geometry.Theory.Segment A B) A\end{verbatim}
\end{theorem}

% === Geometry.Theory.Point ===

\begin{theorem}
  \label{Geometry.Theory.Point.distinct_points_exist}
  \uses{Geometry.Theory.Point, Geometry.Theory.Line, Geometry.Theory.I3}
  \lean{Geometry.Theory.Point.distinct_points_exist}
  \leanok
  For every Point, there is at least one point that isn't that point. 
\begin{verbatim}distinct_points_exist : forall (P : Geometry.Theory.Point), Exists.{1} Geometry.Theory.Point (fun (Q : Geometry.Theory.Point) => Ne.{1} Geometry.Theory.Point P Q)\end{verbatim}
\end{theorem}
